\documentclass[11pt,a4paper]{article}

\usepackage[utf8]{inputenc}
\usepackage[T1]{fontenc}
\usepackage[brazil]{babel}
\usepackage[a4paper,margin=2.5cm]{geometry}
\usepackage{xcolor}
\usepackage{listings}
\usepackage{booktabs}
\usepackage{fancyhdr}
\usepackage[colorlinks=true,linkcolor=blue,urlcolor=blue]{hyperref}

% ---------- Estilo dos blocos de comando ----------
\definecolor{bgcode}{rgb}{0.96,0.96,0.96}
\definecolor{cmt}{rgb}{0.30,0.55,0.30}
\lstdefinestyle{bash}{
  language=bash,
  backgroundcolor=\color{bgcode},
  basicstyle=\ttfamily\small,
  commentstyle=\color{cmt},
  breaklines=true,
  frame=single,
  rulecolor=\color{gray!40},
  columns=fullflexible,
  keepspaces=true,
  showstringspaces=false,
  morecomment=[l]{\#},
  xleftmargin=4pt,
  framexleftmargin=4pt,
}
\lstset{style=bash}

% ---------- Cabeçalho/rodapé ----------
\pagestyle{fancy}
\fancyhf{}
\lhead{Pioneer 3-AT}
\rhead{ROS 2 Humble + Gazebo Harmonic}
\cfoot{\thepage}

\title{\textbf{Pioneer 3-AT em ROS 2 Humble + Gazebo Harmonic}\\[0.3em]
\large Guia de instalação e execução (passo a passo)}
\author{Projeto \texttt{ros-pioneer3at} --- migração ROS 2}
\date{\today}

\begin{document}

\maketitle
\thispagestyle{fancy}

\begin{abstract}
\noindent
Este documento descreve, passo a passo, como instalar e executar a simulação do
robô \textbf{Pioneer 3-AT} em \textbf{ROS 2 Humble} com \textbf{Gazebo Harmonic}
(gz-sim 8). Ao final, você terá o robô em um mundo simulado, controlável pelo
teclado, com odometria, transformadas (TF) e um sensor LiDAR publicados para o
ROS 2 e visualizados no RViz2. Este pacote é uma migração do projeto original em
ROS 1 + Gazebo Classic; o código antigo permanece em \texttt{legacy\_ros1/}.
\end{abstract}

\tableofcontents
\vspace{1em}

% =====================================================================
\section{Pré-requisitos}
% =====================================================================

\begin{itemize}
  \item \textbf{Sistema operacional:} Ubuntu 22.04 (ou base equivalente, ex.\ Zorin OS 17).
  \item \textbf{ROS 2 Humble} (\texttt{ros-humble-desktop}).
  \item \textbf{Gazebo Harmonic} (gz-sim 8.x).
  \item \texttt{git} e \texttt{colcon} para clonar e compilar.
\end{itemize}

\noindent
Confirme as versões antes de começar:

\begin{lstlisting}
source /opt/ros/humble/setup.bash
ros2 --version       # deve indicar Humble
gz sim --version     # deve indicar 8.x (Harmonic)
\end{lstlisting}

% =====================================================================
\section{Instalação}
% =====================================================================

\subsection{Instalar as dependências}

Instale os pacotes de integração ROS$\leftrightarrow$Gazebo e os utilitários
usados pelo projeto:

\begin{lstlisting}
sudo apt update
sudo apt install \
  ros-humble-ros-gz-sim \
  ros-humble-ros-gz-bridge \
  ros-humble-xacro \
  ros-humble-robot-state-publisher \
  ros-humble-joint-state-publisher-gui \
  ros-humble-teleop-twist-keyboard \
  ros-humble-rviz2
\end{lstlisting}

\subsection{Clonar o repositório}

Clone o projeto e selecione a branch da migração ROS 2. Ajuste a URL para o seu
repositório/fork:

\begin{lstlisting}
# escolha uma pasta de trabalho qualquer
cd ~/Documentos
git clone <URL-DO-REPOSITORIO> ros-pioneer3at
cd ros-pioneer3at
git checkout ros2-humble-harmonic-migration
\end{lstlisting}

\subsection{Criar o workspace e compilar}

O ROS 2 compila pacotes dentro de um \emph{workspace} colcon (uma pasta com um
subdiretório \texttt{src/}). Vamos criar um workspace e colocar este pacote
dentro dele por meio de um \emph{symlink} (atalho), de forma que as pastas de
build fiquem fora do repositório:

\begin{lstlisting}
mkdir -p ~/ros2_ws/src
ln -s ~/Documentos/ros-pioneer3at ~/ros2_ws/src/pioneer3at

cd ~/ros2_ws
colcon build --packages-select pioneer3at
\end{lstlisting}

\noindent
Sempre que abrir um terminal novo para usar o projeto, \textbf{carregue o
ambiente} (ROS 2 + workspace):

\begin{lstlisting}
source /opt/ros/humble/setup.bash
source ~/ros2_ws/install/setup.bash
\end{lstlisting}

\noindent
\textit{Dica:} para não repetir isso toda vez, adicione as duas linhas acima ao
final do seu \texttt{\textasciitilde/.bashrc}.

% =====================================================================
\section{Executando a simulação}
% =====================================================================

\subsection{Simulação completa (Gazebo + RViz)}

Sobe, em um único comando, o Gazebo Harmonic com o mundo, o robô, a ponte
ROS$\leftrightarrow$Gazebo e o RViz2:

\begin{lstlisting}
ros2 launch pioneer3at simulation.launch.py rviz:=true
\end{lstlisting}

\noindent
Sem o RViz (apenas Gazebo + ponte):

\begin{lstlisting}
ros2 launch pioneer3at simulation.launch.py
\end{lstlisting}

\noindent
Argumentos disponíveis:

\begin{center}
\begin{tabular}{lll}
\toprule
\textbf{Argumento} & \textbf{Padrão} & \textbf{Descrição} \\
\midrule
\texttt{world}        & \texttt{pioneer\_world.sdf} & arquivo \texttt{.sdf} do mundo \\
\texttt{rviz}         & \texttt{false}              & abre o RViz2 \\
\texttt{use\_sim\_time} & \texttt{true}             & usa o relógio da simulação \\
\bottomrule
\end{tabular}
\end{center}

\subsection{Dirigir o robô pelo teclado (teleop)}

Com a simulação aberta, abra \textbf{outro terminal}, carregue o ambiente
(Seção~2.3) e rode:

\begin{lstlisting}
ros2 run teleop_twist_keyboard teleop_twist_keyboard
\end{lstlisting}

\noindent
Mantenha essa janela em foco e use as teclas indicadas na tela. As principais:

\begin{center}
\begin{tabular}{cl}
\toprule
\textbf{Tecla} & \textbf{Ação} \\
\midrule
\texttt{i} & andar para frente \\
\texttt{,} & andar para tr\'as \\
\texttt{j} & girar para a esquerda \\
\texttt{l} & girar para a direita \\
\texttt{k} & parar \\
\texttt{q} / \texttt{z} & aumentar / diminuir a velocidade \\
\bottomrule
\end{tabular}
\end{center}

\noindent
Os comandos saem no t\'opico \texttt{/cmd\_vel}, s\~ao levados ao Gazebo pela
ponte e movem o rob\^o.

\subsection{Teste rápido sem teclado}

Para verificar o movimento sem o teleop, publique uma velocidade diretamente:

\begin{lstlisting}
ros2 topic pub /cmd_vel geometry_msgs/msg/Twist '{linear: {x: 0.3}}'
\end{lstlisting}

\noindent
O robô deve andar para frente. Pressione \texttt{Ctrl+C} para parar de publicar.

\subsection{Visualizar apenas o modelo (sliders das juntas)}

Para inspecionar a montagem do robô \emph{sem} o Gazebo, usando sliders manuais
para girar as rodas:

\begin{lstlisting}
ros2 launch pioneer3at display.launch.py
\end{lstlisting}

\noindent
Abrem-se o RViz2 (com o modelo) e a janela do
\texttt{joint\_state\_publisher\_gui} (com 4 sliders, um por roda). Arraste os
sliders e confira que cada roda gira em torno do eixo correto e que todos os
links estão conectados na árvore de TF.

% =====================================================================
\section{Tópicos principais}
% =====================================================================

Com a simulação rodando, liste os tópicos com \texttt{ros2 topic list}. Os
principais:

\begin{center}
\begin{tabular}{lll}
\toprule
\textbf{Tópico} & \textbf{Tipo} & \textbf{Sentido} \\
\midrule
\texttt{/cmd\_vel}     & \texttt{geometry\_msgs/msg/Twist}     & ROS $\rightarrow$ Gazebo \\
\texttt{/odom}         & \texttt{nav\_msgs/msg/Odometry}       & Gazebo $\rightarrow$ ROS \\
\texttt{/tf}           & \texttt{tf2\_msgs/msg/TFMessage}      & Gazebo $\rightarrow$ ROS \\
\texttt{/joint\_states} & \texttt{sensor\_msgs/msg/JointState} & Gazebo $\rightarrow$ ROS \\
\texttt{/scan}         & \texttt{sensor\_msgs/msg/LaserScan}   & Gazebo $\rightarrow$ ROS \\
\texttt{/clock}        & \texttt{rosgraph\_msgs/msg/Clock}     & Gazebo $\rightarrow$ ROS \\
\bottomrule
\end{tabular}
\end{center}

\noindent
Exemplos de inspeção:

\begin{lstlisting}
ros2 topic echo /odom          # odometria
ros2 topic echo /scan          # leitura do LiDAR
ros2 topic hz /clock           # taxa do relogio de simulacao
\end{lstlisting}

% =====================================================================
\section{Estrutura do projeto}
% =====================================================================

\begin{lstlisting}
urdf/pioneer3at.urdf.xacro   # descricao do robo (visual + colisao + inercia + plugins)
worlds/pioneer_world.sdf     # mundo do Gazebo Harmonic (chao, luz, obstaculo)
config/ros_gz_bridge.yaml    # mapeamento de topicos ROS <-> Gazebo
config/pioneer3at.rviz       # config do RViz para a simulacao (frame odom)
config/pioneer3at_display.rviz # config do RViz para o display (frame base_link)
launch/
  simulation.launch.py       # bringup completo (gazebo + spawn + bridge + rviz)
  gazebo.launch.py           # so o Gazebo + mundo
  spawn_robot.launch.py      # robot_state_publisher + spawn no gz
  bridge.launch.py           # so a ponte ros_gz_bridge
  display.launch.py          # so RViz com o modelo + sliders (sem Gazebo)
meshes/                      # malhas STL (reaproveitadas do projeto original)
legacy_ros1/                 # codigo ROS 1 / Gazebo Classic original (referencia)
\end{lstlisting}

% =====================================================================
\section{Solução de problemas}
% =====================================================================

\begin{description}
  \item[O robô não aparece no Gazebo / malhas não carregam.]
  Verifique se o ambiente está carregado (Seção~2.3) e recompile. O launch
  define automaticamente \texttt{GZ\_SIM\_RESOURCE\_PATH} para que o Gazebo
  resolva os caminhos \texttt{package://pioneer3at/...}.

  \item[Avisos ``Detected jump back in time'' e o robô pisca no RViz.]
  Ocorre quando o RViz inicia antes do relógio de simulação (\texttt{/clock})
  estar disponível. O \texttt{simulation.launch.py} já mitiga isso iniciando o
  RViz alguns segundos depois do Gazebo. Confirme também que o RViz usa tempo de
  simulação: \texttt{ros2 param get /rviz2 use\_sim\_time} deve retornar
  \texttt{True}.

  \item[\texttt{display.launch.py} abre, mas o robô não aparece.]
  Esse modo usa o frame fixo \texttt{base\_link} (não há \texttt{odom} sem o
  Gazebo). Garanta que está usando a config \texttt{pioneer3at\_display.rviz}
  (já é o padrão do launch).

  \item[Alterei arquivos do pacote e nada mudou.]
  Recompile e recarregue o ambiente:
  \begin{lstlisting}
cd ~/ros2_ws && colcon build --packages-select pioneer3at
source install/setup.bash
  \end{lstlisting}
\end{description}

% =====================================================================
\section{Trabalho futuro}
% =====================================================================

A base de simulação já publica \texttt{/odom}, \texttt{/tf} e \texttt{/scan}, que
são os pré-requisitos para:

\begin{itemize}
  \item \textbf{Navegação autônoma} com o \textbf{Nav2}.
  \item \textbf{Mapeamento (SLAM)} com o \textbf{slam\_toolbox}.
\end{itemize}

\noindent
Esses recursos substituem o antigo \texttt{move\_base} / \texttt{gmapping} /
\texttt{amcl} do projeto em ROS 1.

\end{document}
